\addcontentsline{toc}{section}{СПИСОК ИСПОЛЬЗОВАННЫХ ИСТОЧНИКОВ}

\begin{thebibliography}{53}
	
	\bibitem{lutz}
	Лутц, М. Изучаем Python. Том 1 / М. Лутц. – 5-е изд. – Санкт-Петербург : Питер, 2019. – 832 с. – ISBN 978-5-907144-52-1. – Текст : непосредственный.
	
	\bibitem{ramalho}
	Рамальо, Л. Python. К вершинам мастерства / Л. Рамальо. – Москва : ДМК Пресс, 2017. – 768 с. – ISBN 978-5-97060-384-0. – Текст : непосредственный.
	
	\bibitem{python_docs}
	Python Documentation : официальный сайт. – URL: \url{https://docs.python.org/3/} (дата обращения: 10.01.2026).
	
	\bibitem{pep8}
	PEP 8 – Style Guide for Python Code : официальный сайт. – URL: \url{https://peps.python.org/pep-0008/} (дата обращения: 15.01.2026).
	
	\bibitem{threading}
	Threading --- Thread-based parallelism // Python Documentation : официальный сайт. – URL: \url{https://docs.python.org/3/library/threading.html} (дата обращения: 20.01.2026).
	
	\bibitem{logging}
	Logging HOWTO // Python Documentation : официальный сайт. – URL: \url{https://docs.python.org/3/howto/logging.html} (дата обращения: 25.01.2026).
	
	\bibitem{fastapi_book}
	Любанович, Б. FastAPI: веб-разработка на Python / Б. Любанович. – Санкт-Петербург : Питер, 2024. – 288 с. – ISBN 
	978-601-08-3847-5. – Текст : непосредственный.
	
	\bibitem{fastapi_docs}
	FastAPI Documentation : официальный сайт. – URL: \url{https://fastapi.tiangolo.com/} (дата обращения: 01.02.2026).
	
	\bibitem{pydantic}
	Pydantic Documentation : официальный сайт. – URL: \url{https://docs.pydantic.dev/} (дата обращения: 05.02.2026).
	
	\bibitem{uvicorn}
	Uvicorn Documentation : официальный сайт. – URL: \url{https://www.uvicorn.org/} (дата обращения: 10.02.2026).
	
	\bibitem{starlette}
	Starlette Documentation : официальный сайт. – URL: \url{https://www.starlette.io/} (дата обращения: 14.02.2026).
	
	\bibitem{date}
	Дейт, К. Дж. Введение в системы баз данных / К. Дж. Дейт. – 8-е изд. – Москва : Вильямс, 2005. – 1328 с. – ISBN 978-5-8459-0788-2. – Текст : непосредственный.
	
	\bibitem{connolly}
	Коннолли, Т. Базы данных. Проектирование, реализация и сопровождение. Теория и практика / Т. Коннолли, К. Бегг. – 3-е изд. – Москва : Вильямс, 2003. – 1440 с. – ISBN 978-5-8459-0527-3. – Текст : непосредственный.
	
	\bibitem{groff}
	Грофф, Дж. Р. SQL: полное руководство / Дж. Р. Грофф, П. Н. Вайнберг, Э. Дж. Оппель. – 3-е изд. – Санкт-Петербург : Диалектика, 2020. – 960 с. – ISBN 978-5-907114-26-5. – Текст : непосредственный.
	
	\bibitem{postgres_book}
	Моргунов, Е. П. PostgreSQL. Основы языка SQL / Е. П. Моргунов. – Санкт-Петербург : БХВ-Петербург, 2024. – 336 с. – ISBN 978-5-9775-4022-3. – Текст : непосредственный.
	
	\bibitem{postgres_docs}
	PostgreSQL Documentation : официальный сайт. – URL: \url{https://www.postgresql.org/docs/} (дата обращения: 18.02.2026).
	
	\bibitem{sqlalchemy}
	SQLAlchemy Documentation : официальный сайт. – URL: \url{https://docs.sqlalchemy.org/} (дата обращения: 22.02.2026).
	
	\bibitem{alembic}
	Alembic – Database Migrations for SQLAlchemy : официальный сайт. – URL: \url{https://alembic.sqlalchemy.org/} (дата обращения: 26.02.2026).
	
	\bibitem{psycopg}
	Psycopg – PostgreSQL adapter for Python : официальный сайт. – URL: \url{https://www.psycopg.org/docs/} (дата обращения: 01.03.2026).
	
	\bibitem{python_docx}
	python-docx Documentation : официальный сайт. – URL: \url{https://python-docx.readthedocs.io/} (дата обращения: 05.03.2026).
	
	\bibitem{pywin32}
	pywin32 Documentation // GitHub : официальный сайт. – URL: \url{https://github.com/mhammond/pywin32} (дата обращения: 09.03.2026).
	
	\bibitem{openxml}
	ECMA-376 Office Open XML File Formats : официальный сайт. – URL: \url{https://www.ecma-international.org/publications-and-standards/standards/ecma-376/} (дата обращения: 13.03.2026).
	
	\bibitem{jwt}
	JWT Introduction : официальный сайт. – URL: \url{https://jwt.io/introduction} (дата обращения: 17.03.2026).
	
	\bibitem{oauth}
	The OAuth 2.0 Authorization Framework : официальный сайт. – URL:\url{https://datatracker.ietf.org/doc/html/rfc6749} (дата обращения: 21.03.2026).
	
	\bibitem{martin_clean_code}
	Мартин, Р. Чистый код: создание, анализ и рефакторинг / Р. Мартин. – Санкт-Петербург : Питер, 2016. – 464 с. – ISBN 978-5-496-00487-9. – Текст : непосредственный.
	
	\bibitem{martin_arch}
	Мартин, Р. Чистая архитектура. Искусство разработки программного обеспечения / Р. Мартин. – Санкт-Петербург : Питер, 2018. – 352 с. – ISBN 978-5-4461-0772-8. – Текст : непосредственный.
	
	\bibitem{uml}
	Рамбо, Дж. UML 2.0. Объектно-ориентированное моделирование и разработка / Дж. Рамбо, М. Блаха. – 2-е изд. – Санкт-Петербург : Питер, 2021. – 542 с. – ISBN 978-5-4461-9428-5. – Текст : непосредственный.
	
	\bibitem{uml_spec}
	Unified Modeling Language Specification, Version 2.5.1 : официальный сайт. – URL: \url{https://www.omg.org/spec/UML/2.5.1/PDF} (дата обращения: 25.03.2026).
	
	\bibitem{drawio}
	draw.io Documentation : официальный сайт. – URL: \url{https://www.drawio.com/doc/} (дата обращения: 28.03.2026).
	
	\bibitem{git_book}
	Чакон, С. Git для профессионального программиста / С. Чакон, Б. Штрауб. – Санкт-Петербург : Питер, 2016. – 496 с. – ISBN 978-5-496-01763-3. – Текст : непосредственный.
	
	\bibitem{git_docs}
	Git Documentation : официальный сайт. – URL: \url{https://git-scm.com/doc} (дата обращения: 31.03.2026).
	
	\bibitem{html_css}
	Дэкетт, Д. HTML и CSS. Разработка и создание веб-сайтов / Д. Дэкетт. – Москва : Эксмо, 2014. – 480 с. – ISBN 978-5-699-64193-2. – Текст : непосредственный.
	
	\bibitem{js}
	Дэкетт, Д. JavaScript и jQuery. Интерактивная web-разработка / Д. Дэкетт. – Москва : Эксмо, 2017. – 640 с. – ISBN 978-5-699-80285-2. – Текст : непосредственный.
	
	\bibitem{frein}
	Фрейн, Б. HTML5 и CSS3 и Web 2.0. Разработка современных web-сайтов / Б. Фрейн. – 2-е изд. – Санкт-Петербург : Питер, 2022. – 384 с. – ISBN 978-5-9775-0596-3. – Текст : непосредственный.
	
	\bibitem{mdn_html}
	HTML // MDN Web Docs : официальный сайт. – URL: \url{https://developer.mozilla.org/ru/docs/Web/HTML} (дата обращения: 02.04.2026).
	
	\bibitem{mdn_css}
	CSS // MDN Web Docs : официальный сайт. – URL: \url{https://developer.mozilla.org/ru/docs/Web/CSS} (дата обращения: 04.04.2026).
	
	\bibitem{mdn_js}
	JavaScript // MDN Web Docs : официальный сайт. – URL: \url{https://developer.mozilla.org/ru/docs/Web/JavaScript} (дата обращения: 06.04.2026).
	
	\bibitem{http}
	HTTP // MDN Web Docs : официальный сайт. – URL: \url{https://developer.mozilla.org/ru/docs/Web/HTTP} (дата обращения: 08.04.2026).
	
	\bibitem{fielding}
	Fielding, R. HTTP Semantics / R. Fielding, M. Nottingham, J. Reschke. – RFC 9110. – URL: \url{https://www.rfc-editor.org/rfc/rfc9110.html} (дата обращения: 10.04.2026).
	
	\bibitem{rest}
	Fielding, R. Architectural Styles and the Design of Network-based Software Architectures // University of California : официальный сайт. – URL: \url{https://www.ics.uci.edu/~fielding/pubs/dissertation/rest_arch_style.htm} (дата обращения: 12.04.2026).
	
	\bibitem{pytest}
	Оккен, Б. pytest. Быстрое тестирование приложений на Python / Б. Оккен. – Санкт-Петербург : Питер, 2020. – 320 с. – ISBN  978-1-6805-0240-4. – Текст : непосредственный.
	
	\bibitem{pycharm}
	PyCharm Documentation : официальный сайт. – URL: \url{https://www.jetbrains.com/pycharm/documentation/} (дата обращения: 14.04.2026).
	
	\bibitem{bootstrap}
	Bootstrap Documentation : официальный сайт. – URL: \url{https://getbootstrap.com/docs/5.3/} (дата обращения: 16.04.2026).
	
	\bibitem{docker}
	Docker Documentation : официальный сайт. – URL: \url{https://docs.docker.com/} (дата обращения: 18.04.2026).
	
	\bibitem{nginx}
	NGINX Documentation : официальный сайт. – URL: \url{https://nginx.org/en/docs/} (дата обращения: 20.04.2026).
	
	\bibitem{reportlab}
	ReportLab Documentation : официальный сайт. – URL: \url{https://docs.reportlab.com/} (дата обращения: 22.04.2026).
	
	\bibitem{asyncio}
	asyncio Documentation // Python Documentation : официальный сайт. – URL:\url{https://docs.python.org/3/library/asyncio.html} (дата обращения: 24.04.2026).
	
	\bibitem{knuth}
	Марц Н., Уоррен Дж. Большие данные: принципы и практика построения масштабируемых систем обработки данных в реальном времени / Н. Марц, Дж. Уоррен. – Москва : Диалектика, 2018. – 368 с. – ISBN 978-5-8459-2075-1. – Текст : непосредственный.
	
	\bibitem{orm_book}
	Персиваль, Г.Паттерны разработки на Python: TDD, DDD и событийно-ориентированная архитектура / Г. Персиваль, Б. Грегори. – Санкт-Петербург : Питер, 2025. – 336 с. – ISBN 978-601-12-5028-3. – Текст : непосредственный.
	
	\bibitem{bill_lubanovic}
	Любанович, Б. Простой Python. Современный стиль программирования / Б. Любанович. – Санкт-Петербург : Питер, 2020. – 592 с. – ISBN 978-5-4461-1639-3. – Текст : непосредственный.
	
	\bibitem{daniel_voors}
	Лутц, М. Python. Карманный справочник / М. Лутц. – Москва : Эксмо, 2022. – 318 с. – ISBN 978-5-907114-60-9. – Текст : непосредственный.
	
	\bibitem{python_tutorial}
	Python Tutorial // Python Documentation : официальный сайт. – URL: \url{https://docs.python.org/3/tutorial/} (дата обращения: 26.04.2026).
	
	\bibitem{ramalho_2}
	Рамальо, Л. Python. К вершинам мастерства. Том 2 / Л. Рамальо. – Москва : ДМК Пресс, 2023. – 672 с. – ISBN 978-5-97060-885-2. – Текст : непосредственный.
	
	\bibitem{downey}
	Дауни, А. Б. Think Python. Аналитический подход к программированию / А. Б. Дауни. – 3-е изд. – Санкт-Петербург : Питер, 2024. – 320 с. – ISBN 978-1-4493-3072-9. – Текст : непосредственный.
	
\end{thebibliography}